\documentclass[11pt]{article}
\usepackage[a4paper,margin=1in]{geometry}
\usepackage{amsmath,amssymb,mathtools,bm}
\usepackage{enumitem,booktabs,array}
\usepackage{tcolorbox}
\usepackage{hyperref}
\usepackage[protrusion=true,expansion=false]{microtype}
\hypersetup{colorlinks=true,linkcolor=blue,urlcolor=blue,citecolor=blue}
\setlist[itemize]{topsep=2pt,itemsep=2pt}
\setlist[enumerate]{topsep=2pt,itemsep=2pt}
\newtcolorbox{keybox}{colback=blue!3!white,colframe=blue!40!black,boxrule=0.4pt,arc=1mm}
\newtcolorbox{alertbox}{colback=red!3!white,colframe=red!40!black,boxrule=0.4pt,arc=1mm}
\newtcolorbox{answerbox}{colback=green!3!white,colframe=green!40!black,boxrule=0.4pt,arc=1mm}
\newtcolorbox{drillbox}{colback=yellow!5!white,colframe=orange!60!black,boxrule=0.4pt,arc=1mm}
\newcommand{\EV}{\mathbb{E}}
\newcommand{\Var}{\mathrm{Var}}
\newcommand{\Cov}{\mathrm{Cov}}
\newcommand{\blank}[1]{\underline{\hspace{#1}}}

\title{W10-03: Bates, Kahle \& Stulz (2009, JF) \\
\large ``Why Do U.S.\ Firms Hold So Much More Cash Than They Used To?'' --- Active Recall Workbook}
\author{Banking \& Risk Management --- Exam Prep}
\date{\today}
\begin{document}\maketitle

\begin{keybox}
\textbf{Paper in one sentence.} Average cash-to-assets ratio of US industrials \emph{more than doubled} from $\sim 10\%$ in 1980 to $\sim 23\%$ in 2006; BKS show the rise reflects a \emph{precautionary motive} driven by higher R\&D intensity, higher idiosyncratic cash-flow volatility, and lower inventory/receivables --- structural changes in the US firm, not a simple governance or tax story.

\textbf{Placement.} The canonical paper on the secular cash-holding trend; a natural anchor for precautionary-savings interpretations of corporate liquidity and risk management.
\end{keybox}

\section*{Round 1: Complete Walk-Through}

\subsection*{1.1 The fact}
From Compustat, the average cash-to-assets ratio of US industrial firms rises from $\approx 10.5\%$ in 1980 to $\approx 23.2\%$ in 2006. This holds across size bins; median and mean both rise.

\subsection*{1.2 Framework}
Precautionary cash: firms hold cash to cover future cash-flow shortfalls and to preserve investment. The Opler-Pinkowitz-Stulz-Williamson (1999) cash regression:
\[
\text{Cash/Assets}_{it}=\alpha + \beta_1\cdot\text{R\&D/Assets}_{it} + \beta_2\cdot\text{CFVolatility}_{it} + \beta_3\cdot\text{Size}_{it} + \cdots + \varepsilon_{it}.
\]
BKS estimate the regression within years and use the coefficients to quantify how much of the time-trend is explained by changes in the covariates.

\subsection*{1.3 Headline findings}
\begin{itemize}
\item R\&D intensity rises markedly over the sample, and is strongly positively related to cash.
\item Idiosyncratic cash-flow volatility rises, and is positively related to cash.
\item Inventories and receivables fall (just-in-time technologies), reducing non-cash liquidity; firms substitute cash.
\item Net-debt ratios fall --- cash rises faster than gross debt rises.
\item Time-fixed-effects regression: R\&D, volatility, inventories explain most of the cash increase.
\end{itemize}

\subsection*{1.4 Alternative explanations ruled out}
\begin{itemize}
\item Agency costs: cash rise is not concentrated in poorly governed firms.
\item Taxes/repatriation: effect remains after excluding multinationals.
\item Regulation: secular, not policy-specific.
\end{itemize}

\subsection*{1.5 Implications}
\begin{itemize}
\item US firms became \emph{riskier} at the cash-flow level; cash hoarding is rational precaution.
\item Links to the FSS (1993) framework and Bolton-Chen-Wang dynamic cash models.
\item Important context for Fahlenbrach-Rageth-Stulz (2021, W10-04) on COVID cash value.
\end{itemize}

\subsection*{1.6 Caveats}
\begin{alertbox}
\begin{itemize}
\item Cross-sectional covariates may not fully capture time-series variation.
\item R\&D intensity itself is endogenous to cash availability.
\item Panel could have survivorship bias.
\end{itemize}
\end{alertbox}

\section*{Round 2: Level 1 Fill-in}
\begin{enumerate}
\item Cash ratio: \blank{1cm}\% in 1980 $\to$ \blank{1cm}\% in 2006.
\item Main drivers: R\&D, \blank{2cm}, and \blank{1cm}/receivables.
\item OPSW (\blank{1cm}) cash regression template.
\item Agency story: \blank{1cm} (confirmed/ruled out).
\item Net debt: \blank{1cm} (up/down) over the period.
\end{enumerate}

\section*{Round 3: Level 2 Prompts}
\begin{enumerate}
\item Write the OPSW-style cash regression and identify the key covariates.
\item Explain how BKS decompose the time trend into covariate changes.
\item Why is the agency story insufficient?
\item Link BKS to Fahlenbrach-Rageth-Stulz (2021) on COVID.
\end{enumerate}

\section*{Round 4: Minimal Scaffolding}
From a blank page: (i) secular fact, (ii) framework, (iii) drivers, (iv) ruled-out alternatives, (v) implications.

\section*{Round 5: Blank Page}
Essay: why do US firms hoard cash? The BKS structural explanation.

\vspace{6pt}
\begin{drillbox}
\textbf{Speed Drill.} (1) 1980 vs.\ 2006 cash ratios. (2) Three drivers. (3) Cash regression template paper. (4) Agency story status. (5) Net-debt trend.
\end{drillbox}

\section*{Quick Reference \& Answer Key}
\begin{answerbox}
\begin{itemize}
\item \textbf{Fact.} Cash/assets $\sim 10\%\to\sim 23\%$, 1980--2006.
\item \textbf{Drivers.} R\&D intensity $\uparrow$; idiosyncratic cash-flow volatility $\uparrow$; inventories/receivables $\downarrow$.
\item \textbf{Interpretation.} Precautionary, not agency/tax.
\item \textbf{Net debt.} Actually falls --- cash rises faster than debt.
\end{itemize}
\end{answerbox}

\begin{keybox}
\textbf{30-second exam pitch.} Bates, Kahle \& Stulz (2009) document and explain the secular rise in US corporate cash holdings from about 10\% of assets in 1980 to about 23\% in 2006. Using the Opler--Pinkowitz--Stulz--Williamson cash-determinants regression within years and decomposing the time trend, they show the rise is driven by (i) higher R\&D intensity, (ii) higher idiosyncratic cash-flow volatility, and (iii) declining inventories and receivables (just-in-time management). The increase is a precautionary response to a riskier cash-flow environment, not a governance failure, tax/repatriation artefact, or regulatory effect. Net debt actually declines as cash rises faster than gross debt. The paper reframes US corporate liquidity: firms hoard cash because they have become structurally riskier, in line with Froot--Scharfstein--Stein (1993) precautionary-savings logic.
\end{keybox}

\end{document}
